\section{Method}
\label{sec:method}

\subsection{Cohort and deduplication}
\label{sec:method:cohort}

We use MetalCytoToxDB \citep{krasnov2026metalcytotoxdb} under the source paper's
own cohort-construction rule: keep rows with an IC50 value, keep the five metals
the database covers (Ru, Ir, Rh, Os, Re), apply measurement-level deduplication
on (\texttt{SMILES\_Ligands}, \texttt{Cell\_line}, \texttt{Time(h)}), and apply
the paper's 76-name non-cancer cell-line exclusion. We deliberately do
\emph{not} impose a short-exposure-time window: an earlier internal protocol
restricted rows to $\leq24$\,h, which silently truncated the smaller metals
(the supplement documents the trap). Dropping that window raises the cohort to
$4{,}793$ (Ru), $1{,}281$ (Ir), $334$ (Rh), $333$ (Os) and $250$ (Re) unique
ligands, and is what makes all five metals resolvable. \MEASURED:
\texttt{metrics/by\_metric/w1\_cohort\_expansion\_probe.json}.

Cytotoxicity is binarised from the dark IC50 with a $10\,\mu$M activity
threshold; rows below a $0.01\,\mu$M noise floor are dropped. Ligand-level
duplicates are then collapsed to one row per unique ligand: for every ligand the
label is the majority over its measurements, and exact ties are broken
\emph{toward active} (the predicate is $\text{sum}\times 2 \geq n$). Ties going
to active is the conservative choice for a cytotoxicity endpoint and is the rule
the released code actually implements. Because the same organic ligand can be
measured against several metals, a raw row split leaks ligands across the
train/test boundary; the leave-one-metal-out protocol below removes that leak
explicitly.

\subsection{Evaluation regimes}
\label{sec:method:regimes}

We hold the cohort and the model fixed and vary only the split.

\textbf{Within-metal (ceiling).} For each metal, stratified 5-fold
cross-validation over that metal's unique ligands, ten seeds. This is the
setting the source paper reports, and it is the \emph{ceiling} against which
cross-metal transfer is measured.

\textbf{Cross-metal (LOMO, zero-shot).} For each held-out metal $m$, train on
the unique ligands of the other four metals and test on $m$. The leakage rule is
strict: any training ligand that also occurs in the held-out metal is dropped
from the train pool before fitting, so the test set is ligand-disjoint from
training by construction. This is the \emph{collapse} regime.

\textbf{Few-shot metal adaptation (\fsma).} Starting from the frozen zero-shot
LOMO model, fine-tune on $K$ labelled support ligands drawn from the held-out
metal, then score a query set from the same metal that is disjoint from the
support set. The reported quantity is the \emph{paired} difference between the
adapted query AUC and the \emph{same run's} zero-shot query AUC, so
metal-to-metal difficulty and seed variance cancel.

The three regimes share the cohort and the model; only the split and the
adaptation step differ, which is what lets us read the LOMO drop as a metal
effect rather than a preprocessing artefact.

\subsection{Model}
\label{sec:method:models}

The backbone is the official chemprop~v2 atom message-passing encoder
\citep{chemprop, dmpnn2019} with a single pooled binary-classification head:
$d_h=300$, depth $3$, FFN $1\times300$, dropout $0.0$, trained for $50$ epochs
at batch $64$ and learning rate $10^{-3}$ with class weighting, under the
ligand-disjoint LOMO split. No metal features are supplied
(\texttt{metal\_mode: none}); metal identity enters only through the adaptation
step below. The same architecture and schedule are used for every regime, so the
within-metal ceiling, the zero-shot LOMO column and the \fsma{} curve are
directly comparable. \MEASURED:
\texttt{metrics/by\_metric/pa3\_chemprop\_lomo\_probe\_cohort.json}.

ROC-AUC on held-out ligands, averaged over seeds, is the reported metric. We
also report PR-AUC because the positive rate varies across metals.

\subsection{Few-shot metal adaptation (\fsma)}
\label{sec:method:method}

\fsma{} has four parts.

\textbf{1. Source model.} Pretrain the encoder on the pooled ligands of the
four source metals under the LOMO split, giving a zero-shot scorer
$f_{\theta_0}$ for the held-out metal. This is the model whose query AUC appears
in the cross-metal column.

\textbf{2. Support set.} Draw $K$ labelled ligands from the held-out metal,
disjoint from the query set. By default the draw is label-stratified (each class
contributes proportionally). The support set is the \emph{only} new supervision
\fsma{} uses.

\textbf{3. Adaptation operator.} Freeze the encoder and fine-tune \emph{only
the activity head} for $20$ epochs at learning rate $10^{-3}$. We call this the
head-only scope. The alternative scopes --- full-encoder fine-tuning, and
head plus metal-conditioning modules --- are the subject of the depth ablation
(\S\ref{sec:ablation:scope}); head-only is the default because it is the
scope that keeps the low-$K$ regime non-negative.

\textbf{4. Query evaluation.} Score the disjoint query set with the adapted
model and report $\Delta = \mathrm{AUC}_{\text{adapted}} -
\mathrm{AUC}_{\text{zero-shot}}$ per (metal, seed) cell, then aggregate.

The design is deliberately minimal: one frozen encoder, one trained head, $K$
labels, no architecture change. We sweep $K\in\{5,10,20,40,80\}$ over ten seeds
and all five metals.

\subsection{Controls and ablations}
\label{sec:method:controls}

To separate ``the adaptation operator'' from ``more labels'' and ``more
compute'', every control is capped at the \emph{same optimizer-step count} as
the \fsma{} fine-tune (itself $20$ steps at these $K$): a
\textbf{from-scratch} model trained only on the $K$ support ligands, and a
\textbf{concat-retrain} model trained on the source pool plus the $K$ support
ligands. Both use the identical architecture and the identical support set. A
control that is allowed to train for more steps is not a control, and we show in
\S\ref{sec:ablation:controls} that the un-matched version reverses the
conclusion.

The support-set construction rule is ablated with four samplers: random,
label-stratified (default), Bemis--Murcko scaffold-diverse
\citep{bemis1996murcko}, and activity-scaffold. Adaptation depth is ablated with
three scopes: head-only, head plus metal-conditioning modules, and full-encoder.
The backbone is ablated by re-running the identical protocol with a
fingerprint-plus-XGBoost scorer whose adaptation step is warm-started gradient
boosting (\texttt{xgb\_model} continuation) instead of head fine-tuning.

\subsection{Replication of the published protocol}
\label{sec:method:krasnov}

For the head-to-head against \citet{krasnov2026metalcytotoxdb} we adopt their
criteria rather than our own: measurement-level dedup on
(\texttt{SMILES\_Ligands}, \texttt{Cell\_line}, \texttt{Time(h)}) keeping the
first record, the 76-name non-cancer cell-line exclusion, and their primary
evaluation scheme, a DOI-grouped 10-fold cross-validation in which all rows
sharing a DOI fall in the same fold (DOI-disjointness verified). Our cohort is
larger than theirs --- $15{,}799$ rows / $4{,}793$ complexes / $591$ cell lines,
against $11{,}294$ / $4{,}399$ / $23$ --- so the DOI-split comparison is a
cross-cohort estimate, labelled as such. The like-for-like comparison is the
within-metal axis, where we hold our cohort fixed and compare to the published
within-metal values. \MEASURED:
\texttt{metrics/by\_metric/metalcytotox\_ru\_auc.json::paper\_protocol\_measurement},
\texttt{metalcytotox\_ru\_paper\_lightgbm\_cv.json}.

\subsection{Statistics}
\label{sec:method:stats}

Every claim is a \emph{paired} quantity over (metal $\times$ seed) cells, tested
with a one-sided paired $t$-test as the primary test. A percentile bootstrap is
recorded in the artifacts as a reference only: at $n=5$ seeds it cannot generate
enough distinct resamples and is anti-conservative, so the family correction is
never based on it. Within each \emph{family} --- one $K$-curve, one sampler
sweep, one scope sweep, one control set --- we control the family-wise error
rate with a Holm--Bonferroni step-down correction \citep{holm1979} and report
which cells survive. Per-metal rows are always reported; a pooled row is labelled
Ru-dominated, since Ru is the largest metal. All $p$-values are one-sided in the
direction of the claim. \MEASURED:
\texttt{metrics/by\_metric/fsma\_paper\_stats.json}.

\subsection{Reproducibility}
\label{sec:method:repro}

Every number in this paper traces to a JSON artifact under
\texttt{metrics/by\_metric/} or a run report under \texttt{molmetal/reports/};
the mapping from claim to file is given at each table and figure. Scripts,
configs and cohorts are versioned in the same repository. We follow a strict
MEASURED/PROJECTED discipline: no projected value appears in a results table.
